\documentclass[english,compress]{beamer}
\usepackage{kloeckislides}
\nonstopmode

\usepackage[normalem]{ulem}
\usepackage{pifont}
\usepackage{ifthen}

\setbeamercolor{section in head/foot}{use=structure,bg=structure.fg!25!bg}
\defbeamertemplate*{footline}{split theme}
{%
  \leavevmode%
  \begin{beamercolorbox}[wd=.5\paperwidth,ht=2.5ex,dp=1.125ex]{section in head/foot}%
    \insertsectionnavigationhorizontal{\paperwidth}{\hskip0pt plus1filll}{}%
  \end{beamercolorbox}%
  %\begin{beamercolorbox}[wd=.5\paperwidth,ht=2.5ex,dp=1.125ex]{subsection in head/foot}%
    %\insertsubsectionnavigationhorizontal{.5\paperwidth}{}{\hskip0pt plus1filll}%
  %\end{beamercolorbox}%
}


%\useoutertheme[subsection=false]{miniframes}

\setbeamertemplate{frametitle}[default][center]

\AtBeginDocument{%
  {
    \usebeamercolor{section in head/foot}
  }
  
  \pgfdeclareverticalshading{beamer@headfade}{\paperwidth}
  {%
    color(0cm)=(bg);
    color(1.25cm)=(section in head/foot.bg)%
  }

  \setbeamercolor{section in head/foot}{bg=}
}

\addtoheadtemplate{\pgfuseshading{beamer@headfade}\vskip-1.25cm}{}

\beamertemplatedotitem

\setbeamercolor{section in head/foot}{parent=palette quaternary}
\setbeamercolor{subsection in head/foot}{parent=palette primary}

\setbeamercolor{author in head/foot}{parent=section in head/foot}
\setbeamercolor{title in head/foot}{parent=subsection in head/foot}



\AtBeginSection[] {
  \begin{frame}<beamer>
  \frametitle{Outline}
  \tableofcontents[sectionstyle=show/shaded,subsectionstyle=show/show/hide]
\end{frame}
}
\AtBeginSubsection[] {
  \begin{frame}<beamer>
  \frametitle{Outline}
  \tableofcontents[sectionstyle=show/shaded,subsectionstyle=show/shaded/hide]
\end{frame}
}

\newcommand{\technicality}[2]{%
  {\strut #1\\
    \begin{beamercolorbox}[sep=1mm]{block body}
      #2
    \end{beamercolorbox}
  }%
}

\lstset{
  language=C++,
  rangebeginprefix=//\ ,
  rangeendprefix=//\ ,
}

\def\weblink#1#2{\href{#1}{\color{blue}\underline{#2}}}

\definecolor{fetch}{RGB}{227,110,35}
\definecolor{alu}{RGB}{255,188,24}
\definecolor{context}{RGB}{132,146,175}

\usepackage{keystroke}

\setbeamertemplate{navigation symbols}{}


\begin{document}
% {{{ front matter

\title{High-Performance Scientific Computing\\Lecture 5: More OpenCL, MPI}

\date{MATH-GA 2011 / CSCI-GA 2945 $\cdot$ October 3, 2012}

\frame{\titlepage}

\begin{frame}{Today}
  \tableofcontents[hideallsubsections]
\end{frame}
% }}}
% -----------------------------------------------------------------------------
\begin{frame}{Bits and pieces}
  \begin{itemize}
    \item HW1 grades sent
    \item HW2 graded soon
    \item HW3 due
    \item HW4 out tomorrow
    \item Cuda cluster accounts
    \item Mailing list messages
  \end{itemize}
\end{frame}
% -----------------------------------------------------------------------------
\section{Tool of the day: Git}
% -----------------------------------------------------------------------------
\begin{frame}{Git}
  \begin{center}
  \Huge Demo time
  \end{center}
\end{frame}
% -----------------------------------------------------------------------------
\section{OpenCL: Device Language}
% -----------------------------------------------------------------------------
\begin{frame}{Moar OpenCL!}
  \begin{center}
  \Huge Demo time
  \end{center}
\end{frame}
% -----------------------------------------------------------------------------
% {{{
\begin{frame}{OpenCL Device Language}
  \begin{columns}
    \column{0.7\textwidth}
      OpenCL device language is C99, with these differences:

      \medskip
      \plusball Index getters

      \plusball Memory space qualifiers

      \plusball Vector data types 

      \plusball Many generic (`overloaded') math functions

      \plusball Synchronization

      \minusball Recursion

      \minusball Fine-grained \texttt{malloc()}

      \minusball Function pointers
    \column{0.3\textwidth}
      \includegraphics[width=\textwidth]{opencl-logo.png}
  \end{columns}
\end{frame}

%\input{cuda-cl-dictionary}

\begin{frame}{Address Space Qualifiers}

  \begin{center}
  \begin{tabular}{p{5em}cccp{2.8cm}}
    \hline
    \textbf{Type} & \textbf{Per} & \textbf{``Speed''} \\
    \hline
    private*) & work item & super-fast\\
    local & group & fast \\
    global & grid & kinda slow \\
    \hline
  \end{tabular}

  \bigskip
  *) default, so optional
  \end{center}
  \uncover<2>{
    \begin{tikzpicture} [overlay]
      \node [
        above left=1cm of current page.south east, draw,drop shadow,fill=white,
        inner sep=5mm,thick, text width=0.6\textwidth]
        {
          Should really discuss ``speed'' in terms of
          latency/bandwidth.

          \bigskip
          \emph{Both} decrease with distance from the point of
          execution.
        } ;
    \end{tikzpicture}
  }
\end{frame}

% }}}
% -----------------------------------------------------------------------------
\section{OpenCL: Synchronization}
% -----------------------------------------------------------------------------
% {{{
\begin{nologo}
\begin{frame}{Concurrency and Synchronization}
  \uncover<+->{
    \begin{beamercolorbox}[sep=3mm]{block body}
      GPUs have layers of concurrency.\\
      \hfill Each layer has its synchronization primitives.
    \end{beamercolorbox}
  }
  \bigskip
  \begin{columns}[c]
    \column{0.6\textwidth}
      \uncover<+->{
        \begin{itemize}
          \item Intra-group:\\ 
            \texttt{barrier(\dots)}, \\
            \texttt{mem\_fence(\dots)}\\
            \texttt{\dots} =
            \texttt{CLK\_\{LOCAL,GLOBAL\}\_MEM\_FENCE}
          \item Inter-group:\\ Kernel launch
          \item CPU-GPU:\\ Command queues%, Events
        \end{itemize}
      }
    \column{0.35\textwidth}
      \includegraphics[width=\textwidth]{onion.jpeg}
  \end{columns}

\end{frame}
\end{nologo}
\addimgcredit{Onions: flickr.com/darwinbell \cc}

\input{barrier}
\input{memory-fence}

% -----------------------------------------------------------------------------
\begin{frame}{Synchronization between Groups}
  \begin{block}<+->{Golden Rule:}
    Results of the algorithm must be independent of the order in which
    work groups are executed.
  \end{block}

  \uncover<+->{
  \textbf{Consequences:}
  \begin{itemize}
    \item Work groups may read the same information from global memory.
    \item But: Two work groups may not validly write different things to 
      the same global memory.
    \item Kernel launch serves as
      \begin{itemize}
        \item Global barrier
        \item Global memory fence
      \end{itemize}
  \end{itemize}
  }
\end{frame}
% -----------------------------------------------------------------------------
\begin{frame}{Atomic Operations}
  \uncover<+->{
  Collaborative (inter-block) Global Memory Update:
  \begin{center}
  \begin{tikzpicture}[
    start chain=going right,
    node distance=1cm,
    op/.style={draw,thick,fill=blue!50,on chain},
    every join/.style={},
    ]
    \node [op] (read) {Read} ;
    \node [op] (inc) {Increment} ;
    \node [op] (write) {Write} ;
    \draw [thick,->] (read) -- (inc) coordinate [pos=0.5] (r-i);
    \draw [thick,->] (inc) -- (write) coordinate [pos=0.5] (i-w);
    \uncover<+->{
      \node at (r-i) [arrow box,fill=yellow,
        arrow box arrows={north:0.5cm},below,draw,thick]
        {Interruptible!} ;
    }
    \uncover<+->{
      \node at (i-w) [arrow box,fill=yellow,
        arrow box arrows={north:0.5cm},below,draw,thick]
        {Interruptible!} ;
    }
  \end{tikzpicture}
  \end{center}
  }
  \uncover<+->{
  Atomic Global Memory Update:
  \begin{center}
  \begin{tikzpicture}[
    start chain=going right,
    node distance=1cm,
    op/.style={draw,thick,fill=blue!50,on chain},
    every join/.style={},
    ]
    \node [op] (read) {Read} ;
    \node [op] (inc) {Increment} ;
    \node [op] (write) {Write} ;
    \draw [thick,->] (read) -- (inc) coordinate [pos=0.5] (r-i);
    \draw [thick,->] (inc) -- (write) coordinate [pos=0.5] (i-w);
    \begin{pgfonlayer}{background}
      \node [fit=(read) (write), fill=blue!30,draw,thick] { } ;
    \end{pgfonlayer}
    \uncover<+->{
      \node at (r-i) [arrow box,fill=green!75,
        arrow box arrows={north:0.5cm},below=0.3cm,draw,thick]
        {Protected} ;
    }
    \uncover<+->{
      \node at (i-w) [arrow box,fill=green!75,
        arrow box arrows={north:0.5cm},below=0.3cm,draw,thick]
        {Protected} ;
    }
  \end{tikzpicture}
  \end{center}
  }
  \uncover<+->{
    \textbf{How?}\\
    atomic\_\{add,inc,cmpxchg,\dots\}(int *global, int value);
  }
\end{frame}
% -----------------------------------------------------------------------------
\begin{frame}[fragile]{Atomic: Compare-and-swap}
  \begin{lstlisting}
  int atomic_cmpxchg (__global int *p, int cmp, int val)
  int atomic_cmpxchg (__local int *p, int cmp, int val)
  \end{lstlisting}

  Does:
  \begin{itemize}
    \item Read the 32-bit value (referred to as
    old) stored at location pointed by p.
  \item Compute \texttt{(old == cmp) ? val : old}.
  \item Store result at location pointed by p.
  \item Returns old.
  \end{itemize}
  \uncover<2>{
    \begin{tikzpicture} [overlay]
      \node [above left=1cm of current page.south east, draw,drop shadow,fill=white,
      inner xsep=0.5cm,inner ysep=0.5cm,thick]
        {
          Implement atomic \texttt{float} add?
        } ;
    \end{tikzpicture}
  }
\end{frame}
% }}}
% -----------------------------------------------------------------------------
\section{Intro to MPI}
% -----------------------------------------------------------------------------
% {{{
\begin{frame}{MPI}
  \textbf{M}essage \textbf{P}assing \textbf{I}nterface:
  \begin{center}
    \begin{tikzpicture}
      \node (node1) {
        \includegraphics[height=3cm]{server.jpeg}
      } ;
      \node [right=1cm of node1] (msg) {
        \includegraphics[height=2cm]{envelope.jpeg}
      } ;
      \node [right=1cm of msg] (node2) {
        \includegraphics[height=3cm]{server.jpeg}
      } ;
      \draw [dashed,thick,->] (node1) -- (msg) ;
      \draw [dashed,thick,->] (msg) -- (node2) ;
    \end{tikzpicture}
  \end{center}
\end{frame}
\begin{frame}{MPI}
  \foreach \i in {1,2,...,110}
  {%
    \ifthenelse{\equal{\i}{103}}{%
      \uncover<1-3>{
        \tikz \node (node\i) {\includegraphics[height=0.5cm]{server-small.jpeg}}; %
      }
    }{%
      \tikz \node (node\i) {\includegraphics[height=0.5cm]{server-small.jpeg}}; %
    }
  }
  \uncover<+>{}
  \only<+->{
  \foreach \i in {1,6,...,110}
  {
    \foreach \j in {3,23,...,110}
    {
      \tikz [overlay] \draw [thick,->] (node\i) -- (node\j) ;
    }
  }
  }

  \uncover<+->{
    \begin{tikzpicture} [overlay]
      \node [below right=1cm of current page.north west, draw,drop shadow,fill=white,
      inner xsep=0.5cm,inner ysep=0.5cm,thick]
        {
          Not enough throughput? Just buy more computers$^{*}$
        } ;
    \end{tikzpicture}
  }
  \uncover<+>{}
  \uncover<+->{
    \begin{tikzpicture} [overlay]
      \node [above left=0.5cm of current page.south east, draw,drop shadow,fill=white,
      inner xsep=0.5cm,inner ysep=0.5cm,thick,text width=0.6\textwidth]
        {
          Key questions:
          \begin{itemize}
            \item Who can I send mail to?
            \item How much mail can go through the system (bandwidth)?
            \item How fast does mail arrive (latency)?
            \item Should I wait for the return receipt?
            \only<+>{\item \textbf{Why haven't I heard from the other guys
            yet?}}
          \end{itemize}
        } ;
    \end{tikzpicture}
  }
\end{frame}
\addimgcredit{Server: sxc.hu/Kolobsek}
\addimgcredit{Envelope: sxc.hu/ilco}
% -----------------------------------------------------------------------------
\begin{frame}{MPI}
  \begin{columns}
    \column{0.5\textwidth}
      \begin{center}
      \textbf{MPI 3.0}

      \medskip
      Born September 21, 2012

      \vspace{2cm}
      \footnotesize
      MPI 1.0: June 1994
      \end{center}
    \column{0.5\textwidth}
      \includegraphics[width=\textwidth]{gift-box.png}
  \end{columns}
\end{frame}
\addimgcredit{Gift box: sxc.hu/iprole}
% -----------------------------------------------------------------------------
\begin{frame}{MPI}
  \begin{center}
  \Huge Demo time
  \end{center}
\end{frame}
% -----------------------------------------------------------------------------
\def\hilite<#1>#2{\alt<#1>{\colorbox{blue!30}{#2}}{\colorbox{white}{#2}}}
\begin{frame}{Send definition}
  \textbf{MPI 3.0, Section 3.4:}

  \uncover<+>{}
  \begin{quote}
    \upshape
    [\texttt{MPI\_Send}] is \hilite<+->{blocking}: it does not
    return until the message data and envelope have been safely stored
    away so that the sender is \hilite<+->{free to modify the send buffer}.
  \end{quote}
\end{frame}
% -----------------------------------------------------------------------------
\begin{frame}{Send definition}
  \textbf{MPI 3.0, Section 3.4, more:}

  \uncover<+>{}
  \medskip
  \begin{quote}
    \upshape
    [\texttt{MPI\_Send}] uses the
    \hilite<+->{standard
    \tikz \coordinate (cmode) ;
    communication mode}. In this mode, it is up to
    MPI to decide whether outgoing messages will be buffered.

    \medskip
    MPI
    \tikz \coordinate (may) ;
    \hilite<+->{may} buffer outgoing messages.
    In such a case, the send call may complete before a matching
    receive is invoked. On the other hand, MPI may choose not to
    buffer outgoing messages.
  \end{quote}
  \uncover<+->{
    \begin{tikzpicture} [overlay]
      \node [
        above right=1cm of current page.south west, draw,drop shadow,fill=white,
        inner sep=5mm,thick, text width=0.3\textwidth] (may-descr)
        {
          What is correct behavior?

          \medskip
          Must, should, may (RFC 2119)
        } ;
        \draw [ultra thick,->] (may-descr) -- ($(may)+(2ex,0)$) ;
    \end{tikzpicture}
  }
  \uncover<+->{
    \begin{tikzpicture} [overlay]
      \node [
        above left=5mm of current page.south east, draw,drop shadow,fill=white,
        inner sep=5mm,thick, text width=0.4\textwidth] (cmode-descr)
        {
          Alternative communication modes:
          \begin{itemize}
            \item Buffered
            \item Synchronous
            \item Ready
          \end{itemize}
        } ;
        \draw [ultra thick,->] (cmode-descr) -- (cmode) ;
    \end{tikzpicture}
  }
\end{frame}
% -----------------------------------------------------------------------------
\begin{frame}{Send definition}
  \textbf{MPI 3.0, Section 3.4, yet more:}

  \uncover<+>{}
  \begin{quote}
    \upshape
    A send in standard mode can be \hilite<+->{started} whether or not a matching
    receive has been \hilite<+->{posted}. 

    \medskip
    It may \hilite<+->{complete} before a
    \hilite<+->{matching} receive is posted.

    \medskip
    The standard mode send is \hilite<+->{non-local}: successful
    completion of the send operation may depend on the occurrence of a
    matching receive.
  \end{quote}
\end{frame}
% -----------------------------------------------------------------------------
\begin{frame}{Meta-lesson}
  \begin{center}
    \Large
    Can learn a lot from \emph{how} something is said.
  \end{center}
\end{frame}
% -----------------------------------------------------------------------------
\begin{frame}{Lessons}
  \begin{itemize}
    \item Blocking $\leftrightarrow$ buffers
    \item Communication modes
    \item Operation life cycle
    \item Matching
    \item Non-locality
  \end{itemize}
\end{frame}
% -----------------------------------------------------------------------------
\begin{frame}{Removing the deadlock}
  Two ways laid out:
  \pause
  \begin{itemize}
    \item Use buffered send (brittle!)
    \item Change order (not always easy! Example?)
  \end{itemize}
  \pause
  \bigskip
  Would like a middle ground: 
  \begin{quote}
  ``Just keep the buffer I've got right here!''
  \end{quote}
  But when is it safe to reuse that buffer?
\end{frame}
% -----------------------------------------------------------------------------
\begin{frame}{Non-blocking}
  \textbf{MPI 3.0, Section 3.5:}
  \begin{quote}
    \upshape
    Nonblocking message-passing operations [...] can be used to avoid
    the need for buffering outgoing messages.
  \end{quote}
  Additional Advantage: \textbf{[Sec. 3.7]}
  \begin{quote}
    \upshape
    One can improve performance on many systems by overlapping
    communication and computation.
  \end{quote}
  \uncover<+>{}
  \uncover<+->{
    \begin{tikzpicture} [overlay]
      \node [above left=0.5cm of current page.south east, draw,drop shadow,fill=white,
      inner xsep=0.5cm,inner ysep=0.5cm,thick,text width=0.6\textwidth]
        {
          Nonblocking can be \emph{combined} with
          buffered/ready/synchronous.

          $\rightarrow$ It's not a ``mode''.

          \bigskip
          Nonblocking sends can be matched with blocking receives, and
          vice-versa. \textbf{[3.7]}
        } ;
    \end{tikzpicture}
  }
\end{frame}
% -----------------------------------------------------------------------------
\begin{frame}{MPI}
  \begin{center}
  \Huge Nonblocking demo time
  \end{center}
\end{frame}
% -----------------------------------------------------------------------------
\begin{frame}{Partitioning for neighbor communication}
  \begin{center}
    \includegraphics[height=7cm]{mesh-partition.png}
  \end{center}
  \uncover<+>{}
  \uncover<+->{
    \begin{tikzpicture} [overlay]
      \node [above left=0.5cm of current page.south east, draw,drop shadow,fill=white,
      inner xsep=0.5cm,inner ysep=0.5cm,thick]
        {
          How can I chop up a domain?
        } ;
    \end{tikzpicture}
  }
\end{frame}
% -----------------------------------------------------------------------------
\begin{frame}{MPI}
  \begin{center}
  \Huge Neighbor comm demo time
  \end{center}
\end{frame}
% -----------------------------------------------------------------------------
\begin{frame}{MPI: Ordering}
  \uncover<+->{}
  \textbf{MPI 3.0, Section 3.5:}
  \begin{quote}
    \upshape
    \textbf{Order} Messages are \hilite<+->{non-overtaking}: If a
    sender sends two messages in succession to the same destination,
    and both match the same receive, then this operation cannot
    receive the second message if the first one is still pending.

    \bigskip
    If a receiver posts two receives in succession,
    and both match the same message, then the second receive operation
    cannot be satisfied
    by this message, if the first one is still pending.
  \end{quote}
\end{frame}
% -----------------------------------------------------------------------------
\begin{frame}{MPI: Progress}
  \textbf{MPI 3.0, Section 3.5:}

  \begin{quote}
    \upshape
    \textbf{Progress} If a pair of matching send and receives have
    been initiated on two processes, then at least one of these two
    operations will complete, independently of other actions in the
    system:

    \begin{itemize}
      \item the send operation will complete, unless the receive is
        satisfied by another message, and completes; 
      \item the receive operation will complete, unless the
        message sent is consumed by another matching receive that 
        was posted at the same destination process.
    \end{itemize}
  \end{quote}
\end{frame}
% -----------------------------------------------------------------------------
\begin{frame}[fragile]{MPI: Progress Example}
  \begin{lstlisting}
    if (rank == 0)
    {
      MPI_Bsend(buf1, count, MPI_DOUBLE, 1, tag1, comm)
      MPI_Ssend(buf2, count, MPI_DOUBLE, 1, tag2, comm)
    }
    else if (rank == 1) then
    {
      MPI_Recv(buf1, count, MPI_DOUBLE, 0, tag2, comm, status)
      MPI_Recv(buf2, count, MPI_DOUBLE, 0, tag1, comm, status)
    }
  \end{lstlisting}
\end{frame}



% }}}


% Non-overtaking, order
% Progress

% Using barriers to split 'any' recvs.

% Collectives

% Communicators

% MPE, Jumpshot
% -----------------------------------------------------------------------------



\questionframe{}
\imagecreditslide

\end{document}
% vim: foldmethod=marker
